\section{A concrete Lean formalization route}\label{sec:lean}

\subsection{Repository alignment}

The inspected repository snapshot is
\[
             \code{329161f8da0b319d51ba9f6246c4b19d243f7193}.
\]
The root toolchain is Lean \(4.32.0\); its Lake file
declares the \code{ConciseIPRA} library.
The source namespace is \code{LeanProofs.ConciseIPRA}.
Exact immutable paths and content hashes are recorded
in \code{lean/sources.json}.

The shortest direct route targets the ordinary
iPRA presentation first, then uses its established
translation to the concise scalar presentation.
The inspected interfaces are:
\begin{center}
\begin{tabular}{p{.29\textwidth}p{.62\textwidth}}
\toprule
Purpose & Existing interface\\
\midrule
PR definitions&
\code{PR k}, with \code{zero}, \code{succ}, \code{proj},
\code{comp}, \code{prec}\\
Terms and formulas&
\code{PTm n}, \code{Fm PTm n}, \code{Fm.QF}\\
Object derivability&\code{Prf PRAAx Gamma phi}\\
Restricted induction&\code{PRA.qfInd}\\
Interpretation transport&
\code{pra_theorem_iff}, \code{FromPRA.sigma}\\
\bottomrule
\end{tabular}
\end{center}
The file \code{PR.lean} supplies the language and
induction rule. Interpretation transport comes from
the equivalence in \code{BiInterpretation.lean}
\cite{proveit-pr,proveit-bi}.
The separate DiophantineConsistency axiomatization
was inspected for its explicit distinction between
semantic truth and object-theory provability
\cite{proveit-dioph}. No Busy Beaver or consistency
theorem is needed as a mathematical premise.

\subsection{The three source outlines}

\code{FiniteKernel.lean} defines integer vectors,
residue maps, a finite list of candidate pairs,
a total collision search, and precise input and
output predicates. The decisive host-level target,
still to be proved, says
\[
 \forall c,\ c.\Valid\longrightarrow
 \exists ij,\ 
 \operatorname{collision?}(c)=\operatorname{some}(ij)
 \ \wedge\ c.\operatorname{GoodPair}(ij).
\]
The output relation includes bounds on the two
indices, nonzeroness, and each coordinate of
\(a-b=mz\).

\code{GeometricBridge.lean} supplies symmetry and
midpoint interfaces, with short scripts for the
last logical step and preservation of nonzeroness
by an injective linear map.
The centered-mesh identity is an explicit remaining
algebraic obligation.

\code{ObjectTarget.lean} uses the actual repository
syntax. It packages three function codes of arities
\(1,1,2\), for validation, extraction, and checking,
and forms the quantifier-free formula
\[
                         V(c)=1\longrightarrow G(c,E(c))=1.
\]
It gives only a conditional transport: if a
\code{Prf PRAAx} derivation is provided,
\code{pra_theorem_iff} transports it to the scalar
presentation. The missing derivation is not
postulated as an axiom.

\subsection{Seven completion milestones}

\begin{enumerate}
\item \textbf{Compile against the pinned checkout.}
Resolve elaboration details and inspect dependencies.
The present environment had no Lean or Lake executable;
the files are uncompiled source outlines.
\item \textbf{Implement numerical codecs and PR syntax.}
Give explicit \code{PR.comp} and \code{PR.prec} trees
for signed arithmetic, rational pairs, lists,
finite sums/products, quotient/remainder,
Boolean folds, and bounded search. Prove their
equations in \code{Prf PRAAx}.
\item \textbf{Prove finite pigeonhole in the object theory.}
Formalize the count conservation and no-repeat
invariant of \cref{lem:php}. Convert bounded
checks into bits before using \code{PRA.qfInd}.
\item \textbf{Prove the mesh arithmetic.}
Verify base-\(m\) coding, negative-coordinate
remainders, exact division, and the nonzero
coordinate witness.
\item \textbf{Verify the fixed rational input checker.}
Formalize matrix nonsingularity, slab support tests,
box-interior disjointness, common-denominator
refinement, and extraction of the first \(m^n+1\)
distinct cells.
\item \textbf{Assemble the object derivation.}
Prove the centered-cell identity, basis-map
compatibility, and slab membership. Instantiate
the concrete function symbols and construct
\[
                   \code{Prf PRAAx [] (correctness symbols)}.
\]
The general formula scheme uses the same two
containment instantiations and one midpoint step.
\item \textbf{Transport and audit.}
Apply the repository equivalence, inspect the final
derivation type and dependencies, and record a
reproducible build. Add the finite-volume frontend
and integer-facet enlargement as separate extensions
if the formal target includes
\cref{cor:poly,thm:equality}.
\end{enumerate}

A completed box theorem is already the strict
positive-evidence theorem. A completed formal
theorem starting from a computed polytope volume
must additionally implement the finite-volume
development in \cref{app:volume}.

\subsection{Host-library comparisons}

Mathlib's \code{MeasureTheory} namespace provides
the classical theorem
\begin{lstlisting}
exists_ne_zero_mem_lattice_of_measure_mul_two_pow_lt_measure
\end{lstlisting}
and a compact weak-inequality version
\cite{mathlib-minkowski}. They help compare the
external statement, but do not discharge the
\code{Prf PRAAx} goal.
Similarly, \code{Fintype.exists_ne_map_eq_of_card_lt}
can prove a host-level collision statement
\cite{mathlib-pigeonhole}; the restricted arithmetic
proof still needs the explicit induction audit.

Suggested commands, from the pinned ProveIt checkout
with its dependencies installed, are:
\begin{lstlisting}
lake build ConciseIPRA.BiInterpretation
lake env lean /absolute/path/to/lean/FiniteKernel.lean
lake env lean /absolute/path/to/lean/GeometricBridge.lean
lake env lean /absolute/path/to/lean/ObjectTarget.lean
\end{lstlisting}
These are a completion route, not a build log.
The included Lean directory is not a standalone
Lake project.

\section{Executable artifacts and exact examples}\label{sec:artifacts}

\code{code/minkowski.py} implements the fixed rational
matrix/slab/box format. It checks every input
condition, constructs the mesh, and maintains a
residue table. Arithmetic uses Python integers
and exact \code{fractions.Fraction}; floating-point
input is rejected.
The witness contains the integer vector, a nonzero
coordinate, both cells, their centers, their
physical body points, and the returned lattice
point.

\code{verify_witness} independently checks the
arithmetic relations and body membership.
It verifies enumeration indices by a mixed-radix
rank formula without rerunning the search, and
checks size and volume metadata.
This validates a particular finite witness;
it does not replace the general proof.

\code{code/grid_certificate.py} is an optional
bounded frontend. At a specified mesh and search
radius it enumerates cells, performs exact support
tests, combines adjacent cells into boxes, and
calls the full checker. It trusts no asserted
volume lower bound. Its radius is a search window,
not a verified enclosure: too small a window can
fail to produce a certificate, but cannot certify
an invalid one. The completeness guarantee uses
the independently justified radius and mesh in
\cref{eq:mesh-bound,eq:excess}.

\begin{center}
\begin{tabular}{lrrrrl}
\toprule
Example & \(m\) & Cells & Colors & Examined & Output \(Bz\)\\
\midrule
Square &4&36&16&5&\((0,1)\)\\
Skew lattice &4&36&16&5&\((1,1)\)\\
Diamond &16&264&256&89&\((1,0)\)\\
Three-dimensional cube &4&216&64&5&\((0,0,1)\)\\
\bottomrule
\end{tabular}
\end{center}

For the square, \(K=[-3/2,3/2]^2\), the normalized
box is \([-3/4,3/4]^2\) of content \(9/4\).
The selected cells are \(b=(-3,-3)\) and \(a=(-3,1)\);
their centers are \((-5/8,-5/8)\) and \((-5/8,3/8)\).
Their difference is \((0,1)\), as in
\cref{fig:example}.

For the skew example,
\[
 B=\begin{pmatrix}2&1\\0&1\end{pmatrix},\qquad
 K=\{(x,y):|x-y|\leq3,\ |y|\leq3/2\}.
\]
The same coefficient vector \(z=(0,1)\) maps to
\((1,1)\). The covolume is \(2\) and the body
has area \(18>8\).

For the diamond,
\[
 K=\{(x,y):|x+y|\leq3/2,\ |x-y|\leq3/2\}
   =\{(x,y):|x|+|y|\leq3/2\},
\]
the area is \(9/2>4\). The exact inner grid has
content \(264/256=33/32>1\), an example where a
single large interior rectangle would not suffice.

The rational equality record checks the closed
square \([-1,1]^2\) of area \(4\). Its maximum
integer facet bound is \(1\), so the enlargement
factor is \(3/2\). The strict square witness
\((0,1)\) lies in the original square too.
This extra membership calculation is checked by
\code{make_examples.py}; the equality JSON is
a recorded example, not another input schema for
\code{verify_witness}.

The tests cover exact examples, diamond content,
negative coordinates and quotients, fractional
bases, boxes touching along faces, rejection of
insufficient volume, invalid containment and
singular bases, corrupted witnesses, and
noninteger metadata. Nine test methods pass;
the log is included. These are implementation
checks, not a proof of the unbounded theorem by
testing.

From the package root:
\begin{lstlisting}
python3 code/make_examples.py
python3 code/test_minkowski.py
python3 code/minkowski.py examples/diamond.json \
  --verify examples/diamond.witness.json
\end{lstlisting}
Run \code{make} in \code{paper/} to rebuild the
article. The figure PDF is included, so LaTeX
does not require plotting libraries.
Regenerating the figure uses
\code{code/make_figure.py} and Matplotlib.
